\section{Fixed-sign theorem for non-virtual denominators}
\label{app:fixed-sign}

Let $k_a,k_b$ be future-directed incoming null momenta and let
$p_1,\ldots,p_n$ be future-directed outgoing null momenta, with
\begin{equation}
 k_a+k_b=\sum_{i=1}^{n}p_i,
 \qquad s=2k_a\cdot k_b>0.
 \label{eq:sign-theorem-kinematics}
\end{equation}
For real $a,b,c_i$, define
\begin{equation}
 Q=ak_a+bk_b+\sum_{i=1}^{n}c_ip_i,
 \qquad
 \mu_{ij}=(a+c_i)(b+c_j),
 \quad i\ne j.
 \label{eq:sign-theorem-Q}
\end{equation}
Then, on the closure of the physical massless phase space,
\begin{equation}
 \min_{i\ne j}\mu_{ij}
 \leq\frac{Q^2}{s}\leq
 \max_{i\ne j}\mu_{ij},
 \label{eq:fixed-sign-bound}
\end{equation}
and both bounds are attained as boundary limits.  Therefore $Q^2\geq0$ if
all $\mu_{ij}\geq0$, and $Q^2\leq0$ if all
$\mu_{ij}\leq0$.  Equality on a soft or collinear boundary is allowed.  If
the set $\{\mu_{ij}\}$ contains both signs, the theorem does not certify the
denominator.

To prove Eq.~\eqref{eq:fixed-sign-bound}, introduce
\begin{equation}
 x_i=\frac{2k_b\cdot p_i}{s},
 \qquad
 y_i=\frac{2k_a\cdot p_i}{s},
 \qquad
 z_{ij}=\frac{2p_i\cdot p_j}{s}=z_{ji}.
 \label{eq:sign-edge-data}
\end{equation}
These quantities are nonnegative and obey
\begin{equation}
 \sum_i x_i=\sum_i y_i=1,
 \qquad
 x_i+y_i=\sum_{j\ne i}z_{ij},
 \qquad
 \sum_{i<j}z_{ij}=1.
 \label{eq:sign-edge-relations}
\end{equation}
For a subset $A\subseteq\{1,\ldots,n\}$, let
$K_A=\sum_{i\in A}p_i$.  Since $k_b-K_A$ is the difference between one
incoming beam and a subset of the outgoing momenta,
\begin{equation}
 (k_b-K_A)^2\leq0
 \quad\Longrightarrow\quad
 \sum_{i<j\in A}z_{ij}\leq\sum_{i\in A}x_i.
 \label{eq:subset-inequality}
\end{equation}
The analogous inequality holds with $x_i$ replaced by $y_i$.  These are
the max-flow/min-cut conditions for nonnegative numbers
$\lambda_{ij}$, $i\ne j$, satisfying
\begin{equation}
 \lambda_{ij}+\lambda_{ji}=z_{ij},
 \qquad
 \sum_{j\ne i}\lambda_{ij}=x_i,
 \qquad
 \sum_{i\ne j}\lambda_{ij}=y_j,
 \qquad
 \sum_{i\ne j}\lambda_{ij}=1.
 \label{eq:edge-orientation}
\end{equation}
Expanding $Q^2$ and using these marginals gives
\begin{equation}
 \frac{Q^2}{s}
 =\sum_{i\ne j}\lambda_{ij}(a+c_i)(b+c_j)
 =\sum_{i\ne j}\lambda_{ij}\mu_{ij}.
 \label{eq:sign-convex-combination}
\end{equation}
Thus $Q^2/s$ is a convex combination of the $\mu_{ij}$, proving the
bounds.  The boundary in which $p_i\to k_a$, $p_j\to k_b$, and all other
outgoing momenta become soft gives $Q^2/s\to\mu_{ij}$.  Every extremal value
is therefore a physical boundary limit.
